% GENERATED FILE. Edit canonical lessons, then run npm run generate:books.
% canonical-source-hash: fnv1a64:c20188bc02c60a43
% canonical-lessons: TE-C116-sambaru, TE-C116-rasam, TE-C116-guddu, TE-C116-cepa, TE-C116-mamsam

\chapter{Sambar, Rasam, Egg, Fish, Meat}
\label{ch:te-sambar-rasam-egg-fish-meat}
\hlchaptermodality{116}

\begin{chapteropening}
\textbf{By the end of this chapter:} \emph{I can say sambar, rasam, an egg, a fish and meat.}

Complete the last lesson of chapter 116: i can say sambar, rasam, an egg, a fish and meat.
\end{chapteropening}

\section[\texorpdfstring{\te{సాంబారు} (sāmbāru)}{సాంబారు (sāmbāru)}]{\te{సాంబారు} (sāmbāru) — sambar}

\label{lesson:TE-C116-sambaru}



\begin{quote}
Before the new one: say the Telugu for a curry, then the Telugu for dal.
\end{quote}

\subsection*{You'll want to know: \te{సాంబారు}}
\textbf{\te{సాంబారు}} — \emph{sāmbāru} — \textquotedblleft{}sambar\textquotedblright{}.

\begin{grammarlens}[title={What you've built}]
One more word for this chapter.
\end{grammarlens}

\subsection*{Your turn}
\begin{itemize}
\raggedright
  \item \emph{Say it:} \emph{sāmbāru}
  \item \emph{Say it:} \emph{sāmbāru}, once more
  \item \emph{Say it:} say \emph{sāmbāru}
  \item \emph{Recall:} say the Telugu for a curry, then the Telugu for dal, then say \emph{sāmbāru} again
\end{itemize}

\begin{tcolorbox}[breakable,colback=teal!4,colframe=teal!35!black,title={Before you move on}]
What does \te{సాంబారు} mean? (\textquotedblleft{}Sambar\textquotedblright{}.) Say it once more.
\end{tcolorbox}
\section[\texorpdfstring{\te{రసం} (rasaṁ)}{రసం (rasaṁ)}]{\te{రసం} (rasaṁ) — rasam}

\label{lesson:TE-C116-rasam}



\begin{quote}
Before the new one: say the Telugu for dal, then the Telugu for sambar.
\end{quote}

\subsection*{You'll want to know: \te{రసం}}
\textbf{\te{రసం}} — \emph{rasaṁ} — \textquotedblleft{}rasam\textquotedblright{}.

\begin{grammarlens}[title={What you've built}]
One more word for this chapter.
\end{grammarlens}

\subsection*{Your turn}
\begin{itemize}
\raggedright
  \item \emph{Say it:} \emph{rasaṁ}
  \item \emph{Say it:} \emph{rasaṁ}, once more
  \item \emph{Say it:} say \emph{rasaṁ}
  \item \emph{Recall:} say the Telugu for dal, then the Telugu for sambar, then say \emph{rasaṁ} again
\end{itemize}

\begin{tcolorbox}[breakable,colback=teal!4,colframe=teal!35!black,title={Before you move on}]
What does \te{రసం} mean? (\textquotedblleft{}Rasam\textquotedblright{}.) Say it once more.
\end{tcolorbox}
\section[\texorpdfstring{\te{గుడ్డు} (guḍḍu)}{గుడ్డు (guḍḍu)}]{\te{గుడ్డు} (guḍḍu) — an egg}

\label{lesson:TE-C116-guddu}



\begin{quote}
Before the new one: say the Telugu for sambar, then the Telugu for rasam.
\end{quote}

\subsection*{You'll want to know: \te{గుడ్డు}}
\textbf{\te{గుడ్డు}} — \emph{guḍḍu} — \textquotedblleft{}an egg\textquotedblright{}.

\begin{grammarlens}[title={What you've built}]
One more word for this chapter.
\end{grammarlens}

\subsection*{Your turn}
\begin{itemize}
\raggedright
  \item \emph{Say it:} \emph{guḍḍu}
  \item \emph{Say it:} \emph{guḍḍu}, once more
  \item \emph{Say it:} say \emph{guḍḍu}
  \item \emph{Recall:} say the Telugu for sambar, then the Telugu for rasam, then say \emph{guḍḍu} again
\end{itemize}

\begin{tcolorbox}[breakable,colback=teal!4,colframe=teal!35!black,title={Before you move on}]
What does \te{గుడ్డు} mean? (\textquotedblleft{}An egg\textquotedblright{}.) Say it once more.
\end{tcolorbox}
\section[\texorpdfstring{\te{చేప} (cēpa)}{చేప (cēpa)}]{\te{చేప} (cēpa) — a fish}

\label{lesson:TE-C116-cepa}



\begin{quote}
Before the new one: say the Telugu for rasam, then the Telugu for an egg.
\end{quote}

\subsection*{You'll want to know: \te{చేప}}
\textbf{\te{చేప}} — \emph{cēpa} — \textquotedblleft{}a fish\textquotedblright{}.

\begin{grammarlens}[title={What you've built}]
One more word for this chapter.
\end{grammarlens}

\subsection*{Your turn}
\begin{itemize}
\raggedright
  \item \emph{Say it:} \emph{cēpa}
  \item \emph{Say it:} \emph{cēpa}, once more
  \item \emph{Say it:} say \emph{cēpa}
  \item \emph{Recall:} say the Telugu for rasam, then the Telugu for an egg, then say \emph{cēpa} again
\end{itemize}

\begin{tcolorbox}[breakable,colback=teal!4,colframe=teal!35!black,title={Before you move on}]
What does \te{చేప} mean? (\textquotedblleft{}A fish\textquotedblright{}.) Say it once more.
\end{tcolorbox}
\section[\texorpdfstring{\te{మాంసం} (māṁsaṁ)}{మాంసం (māṁsaṁ)}]{\te{మాంసం} (māṁsaṁ) — meat}

\label{lesson:TE-C116-mamsam}



\begin{quote}
Before the new one: say the Telugu for an egg, then the Telugu for a fish.
\end{quote}

\subsection*{You'll want to know: \te{మాంసం}}
\textbf{\te{మాంసం}} — \emph{māṁsaṁ} — \textquotedblleft{}meat\textquotedblright{}.

\begin{grammarlens}[title={What you've built}]
One more word for this chapter.
\end{grammarlens}

\subsection*{Your turn}
\begin{itemize}
\raggedright
  \item \emph{Say it:} \emph{māṁsaṁ}
  \item \emph{Say it:} \emph{māṁsaṁ}, once more
  \item \emph{Say it:} say \emph{māṁsaṁ}
  \item \emph{Recall:} say the Telugu for an egg, then the Telugu for a fish, then say \emph{māṁsaṁ} again
\end{itemize}

\begin{tcolorbox}[breakable,colback=teal!4,colframe=teal!35!black,title={Before you move on}]
What does \te{మాంసం} mean? (\textquotedblleft{}Meat\textquotedblright{}.) Say it once more.
\end{tcolorbox}
